% !TeX program = xelatex
% Evidence-aligned manuscript draft. DCVR-v2 is a proposal, not an evaluated method.
\documentclass[journal]{IEEEtran}
\usepackage{fontspec}
\setmainfont{TeX Gyre Termes}
\usepackage{xeCJK}
\setCJKmainfont[BoldFont=SimHei]{SimSun}
\setCJKsansfont{Microsoft YaHei}
\setCJKmonofont{SimSun}
\usepackage{cite}
\usepackage{amsmath,amssymb,amsfonts}
\usepackage{graphicx,booktabs,array,multirow}
\usepackage{balance}
\usepackage{xcolor}
\usepackage[hidelinks]{hyperref}
\renewcommand{\abstractname}{摘\quad 要}
\renewcommand{\figurename}{图}
\renewcommand{\tablename}{表}
\renewcommand{\IEEEkeywordsname}{关键词}
\newcommand{\R}{\mathrm{R}}
\newcommand{\T}{\mathrm{T}}
\newcommand{\softmax}{\operatorname{softmax}}
\newcommand{\sigmoid}{\operatorname{sigmoid}}
\newcommand{\pending}{\textit{待验证}}
\begin{document}

\title{面向模态退化的 RGB-T 跟踪前景条件协同 Value 路由机制}
\author{\textit{作者与单位待补充}}
\maketitle

\begin{abstract}
RGB-T 单目标跟踪利用可见光与热红外的互补信息，在复杂环境中持续定位目标。由于照明、热背景与遮挡等因素会造成两种模态及其局部区域的可靠性差异，跨模态融合不仅需要建立信息交互，还需要判断消息应由哪个方向传递、由哪些模板位置提供。现有模板多头自注意力依据查询—键相关性聚合跨模态 Value，但相关性本身并不刻画来源 token 对目标定位的贡献。本文围绕这一接口提出退化感知协同 Value 路由，并将其与原生 MSHA 构成耦合融合模块。该模块利用模板—搜索匹配分布、跨模态一致性与首帧目标前景先验，联合分配双向消息预算和来源 token 权重，同时保留原有 QK 匹配与后续定位路径。与整模态质量加权及搜索侧交互方法不同，本文将研究边界限定在模板 MSHA 的跨模态 Value 路径，聚焦前景条件下的消息来源与贡献分配。实验部分整理已有先导结果和基线参照，并给出机制诊断及闭环评估方案；DCVR-v2 的跟踪性能将在匹配评测完成后报告。
\end{abstract}

\begin{IEEEkeywords}
RGB-T 单目标跟踪，跨模态融合，前景条件路由，Value 消息，退化鲁棒性。
\end{IEEEkeywords}

\section{引言}
单目标跟踪根据首帧给定的目标框，在后续视频中持续估计目标位置。安防监控中的低照度、强光、局部遮挡、运动模糊与复杂热背景会改变目标外观及其与背景的可分性。可见光（RGB）提供丰富纹理，但易受照明变化影响；热红外（TIR）具有不同的成像特性，也会受到热交叉和温差不足的干扰。基准研究推动了 RGB-T 跟踪从小规模验证走向多场景评测，并进一步揭示了模态有效性随场景变化的现象\cite{li2019benchmark,li2022lasher,tang2025mvrgbt}。因此，融合的关键不只是组合两路特征，而是使信息交互适应当前目标及其周围区域的实际状态。

近年来，Transformer 跟踪器通过模板与搜索区域之间的注意力交互组织目标信息。模板桥接类方法进一步以模板特征连接双模态搜索表征，在统一架构中建立跨模态交互\cite{hui2023tbsi,li2025tbsi}。这一设计提升了模板信息的作用，但也使模板 Value 成为跨模态消息的来源：模板网格同时包含目标前景、边界与背景，注意力中的 QK 相关性表达查询与键的匹配关系，并不等同于来源位置对目标定位的效用。当某一模态受退化影响，或背景产生强响应时，未经区分的 Value 消息仍可能进入另一模态的目标表示。

由此，本文聚焦于模板交互中的一个具体问题：
\emph{当 RGB 与 TIR 的当前可靠性不对称时，如何在保留原有匹配结构的前提下，依据目标前景和当前匹配证据，协同决定跨模态 Value 消息的传递方向及其来源 token 贡献？}
该问题不同于重新设计搜索区域交互、更新策略或全局模态选择。本文将操作范围限定在原生 MSHA 的双向跨模态 Value 支路：以模板—搜索匹配分布、跨模态一致性和首帧框生成的软前景权重形成路由依据，分别调节方向级消息预算与来源 token 权重；原有 QK 匹配、同模态信息及后续定位路径保持不变。由方向和 token 两个尺度组成的退化感知协同 Value 路由（DCVR-v2）与 MSHA 共同构成本文研究的耦合融合模块。

这一问题定位来自对现有跟踪失败链和融合接口的分析。部分失跟序列在目标仍处于搜索裁剪范围时已出现局部定位偏差，随后误差经状态反馈累积；这表明仅观察最终框是否成功，难以区分表征交互与状态更新各自的作用。已有先导实验进一步表明，跨模态 Value 路径可以被独立调制，但整路缩放不足以表达前景与背景 token 的差异，且局部机制变化尚不能直接推导为闭环跟踪收益。基线参照、先导数据及其证据边界集中在实验部分报告，避免将不同设置下的记录混为同一性能结论。

本文的贡献概括如下：
\begin{itemize}
  \item 将 RGB-T 模板融合中的研究问题具体化为前景条件下的跨模态 Value 消息分配，指出 QK 匹配关系与来源 token 的目标定位效用并不等价。
  \item 提出 DCVR-v2，并与原生 MSHA 构成耦合融合机制。该机制利用当前匹配分布、跨模态一致性和模板前景先验，协同调节双向消息预算及来源 token 贡献。
  \item 明确方法与实验的验证边界：通过固定状态诊断、受控模态退化和匹配闭环评测分别检验路由行为、定位表征及最终跟踪效果，并整理已有实验资产作为后续比较依据。
\end{itemize}

\begin{figure}[!t]
\centering
\fbox{\parbox[c][3.0cm][c]{0.92\columnwidth}{\centering
\textbf{图件预留：问题动机与方法边界}\par
建议展示目标前景/背景模板 token、模态不对称退化以及原生 MSHA 跨模态 Value 路径，并标明本文仅调整的消息分配接口。\par
\textit{正式图需替换为与实际实现一致的结构图。}}}
\caption{RGB-T 退化条件下模板跨模态 Value 消息分配问题及本文方法范围。}
\label{fig:motivation}
\end{figure}

\section{相关工作与问题定位}
\subsection{RGB-T 跟踪与模板交互}
RGB-T 跟踪研究从基准构建与双模态特征融合逐步发展到基于 Transformer 的跨区域交互。早期基准工作建立了 RGB-T 目标跟踪任务及评测协议，LasHeR 进一步扩展了场景规模和属性覆盖，使模态质量变化能够在多样环境下进行分析\cite{li2019benchmark,li2022lasher}。近期基准研究从模态有效性角度重新审视 RGB-T 跟踪，强调不同序列中两种传感器提供的信息并非恒定\cite{tang2025mvrgbt}。在模型结构上，TBSI 通过模板桥接组织 RGB 与 TIR 搜索特征的交互，后续工作延续并拓展了模板保留与更新机制\cite{hui2023tbsi,li2025tbsi}；跨模态蒸馏路线则通过模态间知识迁移促进双模态表示学习\cite{zhang2023crossdistill}。这些研究说明模板既承担目标先验，也参与跨模态表征形成。

本文与模板桥接方法共享模板交互这一结构背景，但研究对象并非模板更新或搜索侧的交互位置。本文保留原生注意力的查询—键匹配及同模态分支，只考察模板来源的跨模态 Value：当模板网格同时含有前景与背景时，如何在方向层面控制消息预算，并在 token 层面区分其相对贡献。因而，本文关注的是既有模板交互内部的消息路由，而非另设一条搜索—模板信息通路。

\subsection{质量感知融合与路由决策}
针对 RGB-T 模态质量变化，已有方法从质量评估、时空建模、模板不确定性和候选筛选等角度调整特征利用方式\cite{ding2025qstnet,ding2025tumnet,zhang2025btmtrack,tang2025mvrgbt}。另有工作研究测试时适配、时序上下文或 token 级路由，以改善跨场景的模态利用与计算分配\cite{shao2025pura,liu2026mosstrack,ding2026routing,guo2025dreamtrack,ni2026etmtrack}。这些路线分别作用于质量估计、时间信息、候选筛选、模型适配或整体 token 计算路径，为动态多模态融合提供了不同机制。

不确定性估计与校准研究则讨论预测分散度、模型间差异及置信度和准确率的一致性\cite{kendall2017uncertainties,lakshminarayanan2017deep,guo2017calibration}。本文借鉴“依据当前证据调节信息使用”的一般思路，但不将 DCVR-v2 定义为通用不确定性估计器，也不声称由匹配熵直接得到经过校准的概率。其输出是针对模板 Value 路径的相对路由权重，目标是控制信息进入另一模态表示的方向与来源位置。

\subsection{问题定义与研究边界}
本文研究模板空间中前景与背景并存、且 RGB/TIR 当前证据不对称时的跨模态 Value 分配。给定原生 MSHA 的双向注意力矩阵及两模态模板 Value，DCVR-v2 从模板—搜索匹配分布、跨模态一致性和首帧目标框诱导的软前景权重中构造路由状态，并联合产生方向级预算与 token 级相对权重。该设计与全局模态质量加权的区别在于：决策落在具体模板消息及其传递方向；与搜索 token 路由的区别在于：它不改变接收端的搜索交互结构；与预测不确定性估计的区别在于：权重仅服务于当前 Value 消息分配，不作为校准置信度解释。

研究范围限定为 RGB-T 单目标跟踪中原生 MSHA 的跨模态 Value 支路。本文不修改 QK 注意力、同模态 Value、搜索区域交互、框解码、模板更新、状态提交及重检测机制，也不把单模态同时失效作为该模块可独立解决的问题。方法是否有效需分别由路由行为、同一裁剪状态下的定位表征变化及闭环跟踪结果支持；因此，机制可解释性不能替代匹配协议下的最终性能评估。

\section{方法}
\subsection{整体框架}
图~\ref{fig:framework} 给出 DCVR-v2 的计算路径。模块接收原生模板 MSHA 的双模态 token 和首帧目标框生成的软前景权重，先从模板—搜索匹配分布提取 token 级证据，再汇聚得到两模态状态；方向分支分配两条跨模态通路的相对预算，token 分支在来源模板内部重排消息贡献。最终仅对跨模态 Value 项进行调制，QK 注意力、同模态消息及后续定位与状态更新保持不变。该图目前留作正式方法图版式，最终结构需与代码实现逐项核对后替换。

\begin{figure*}[!t]
\centering
\fbox{\parbox[c][4.2cm][c]{0.94\textwidth}{\centering
\textbf{图件预留：DCVR-v2 总体结构图}\par
建议流程：双模态模板/搜索 token 与首帧前景权重\; $\rightarrow$\; 匹配分布证据\; $\rightarrow$\; 前景条件模态状态\; $\rightarrow$\; 有向预算与 token 权重\; $\rightarrow$\; MSHA 跨模态 Value 更新\; $\rightarrow$\; 原有定位头。\par
\textit{当前为版式占位；待实现合同确认后绘制，不表示算子已完成验证。}}}
\caption{前景条件协同 Value 路由的拟议计算流程。}
\label{fig:framework}
\end{figure*}

\subsection{跟踪器与符号}
令 $m\in\{\R,\T\}$ 表示模态，$T_m\in\mathbb{R}^{M\times d}$ 和 $S_m\in\mathbb{R}^{N\times d}$ 分别为同模态模板与搜索 token。首帧标注框映射为模板软前景权重 $w_i^m\in[0,1]$。原生 MSHA 使用查询—键分数生成注意力，并在模板 token 间聚合 Value；DCVR-v2 保留该匹配分数，只调节跨模态来源 Value 的权重。对接收模态 $n$，其来自模态 $m$ 的跨模态消息记为 $M_{n\leftarrow m}$，原生同模态消息保持不变。

为明确后续公式中的变量含义，表~\ref{tab:notation} 汇总主要符号。表中描述的是拟议算子定义，不意味着相关量已在实验中验证为可靠性概率。
\begin{table}[!t]
\centering
\caption{DCVR-v2 的主要符号}
\label{tab:notation}
\small
\begin{tabular}{@{}p{1.8cm}p{4.9cm}@{}}
\toprule
符号 & 含义 \\
\midrule
$T_m,S_m$ & 模态 $m$ 的模板 token 与搜索 token \\
$w_i^m$ & 首帧框映射得到的第 $i$ 个模板 token 软前景权重 \\
$e_i^m$ & 匹配强度、分布熵、竞争差及跨模态一致性组成的证据向量 \\
$q_m$ & 由前景 token 汇聚的模态状态 \\
$b_{m\rightarrow n}$ & $m$ 到 $n$ 的跨模态方向预算 \\
$\hat g_i^{m\rightarrow n}$ & 来源 token 的前景归一化相对权重 \\
$A,V$ & 跨模态注意力子矩阵及来源 Value \\
\bottomrule
\end{tabular}
\end{table}

\subsection{模板条件的匹配分布证据}
对来源模板 token $i$，先计算它与同模态搜索 token $j$ 的归一化相似度：
\begin{align}
 s_{ij}^{m}&=\frac{\langle \operatorname{norm}(W_T^mT_{m,i}),
                   \operatorname{norm}(W_S^mS_{m,j})\rangle}{\tau_s},\\
 p_{ij}^{m}&=\softmax_j(s_{ij}^{m}).
\end{align}
从分布而非单一最大值提取四类证据：
\begin{align}
 p_i^m &= \sum_j p_{ij}^m s_{ij}^m, \\
 h_i^m &= -\frac{1}{\log N}\sum_j p_{ij}^m\log(p_{ij}^m+\epsilon), \\
 c_i^m &= \max_{j\in\mathcal{N}_i}s_{ij}^m-\max_{j\notin\mathcal{N}_i}s_{ij}^m, \\
 a_i^m &= \operatorname{cos}(T_{m,i},T_{\bar m,\pi(i)}).
\end{align}
其中 $j_i^*=\arg\max_j s_{ij}^m$，$\mathcal{N}_i$ 是以 $j_i^*$ 为中心、尺度由首帧目标框确定的搜索邻域，$\pi(i)$ 是两模态模板 token 的空间对应索引。$p_i^m$ 表示匹配强度，$h_i^m$ 表示分布不确定性，$c_i^m$ 比较局部峰与非邻域峰，$a_i^m$ 表示对应模板 token 的跨模态一致性。它们只是路由输入，不是校准概率，也不天然等于目标定位效用。令 $e_i^m=[p_i^m,h_i^m,c_i^m,a_i^m]$。

\subsection{前景条件模态状态}
利用首帧框得到的软权重计算模态级状态：
\begin{equation}
 q_m=\frac{\sum_i w_i^m\,\phi(e_i^m)}
           {\sum_i w_i^m+\epsilon},\qquad
 \bar a_{\R\T}=\frac{\sum_i w_i^\R w_{\pi(i)}^\T a_i^\R}
 {\sum_i w_i^\R w_{\pi(i)}^\T+\epsilon}.
\end{equation}
其中 $\phi$ 是共享的小型映射。背景模板 token 不直接支配模态级状态；软边界可降低目标框与 patch 网格不对齐造成的硬切分误差。前景 mask 的生成与使用必须在训练和推理中一致。缺失 mask 时不得静默使用全一权重；应在数据入口显式构造，或将该样本标为协议失败。

\subsection{方向预算与 token 级分配}
对每个有向跨模态通路，方向网络读取两模态状态、跨模态一致性及状态差：
\begin{equation}
 d_{m\rightarrow n}=f_{\mathrm{dir}}([q_m,q_n,\bar a_{mn},q_m-q_n]).
\end{equation}
对两个相反的跨模态方向进行归一化，得到有界方向预算：
\begin{equation}
 \begin{bmatrix}b_{\R\rightarrow\T}\\b_{\T\rightarrow\R}\end{bmatrix}
 =2\,\softmax\!\left(\begin{bmatrix}d_{\R\rightarrow\T}\\
 d_{\T\rightarrow\R}\end{bmatrix}\right),\qquad
 b_{\R\rightarrow\T}+b_{\T\rightarrow\R}=2.
\end{equation}
因而一个方向可相对增强时，另一个方向的预算相应减少；两路平均预算为 1。该预算只作用于跨模态支路，不缩放同模态 Value。

来源 token 网络逐 token 读取本模态与接收模态证据：
\begin{equation}
 z_i^{m\rightarrow n}=f_{\mathrm{tok}}([e_i^m,e_{\pi(i)}^n,w_i^m,a_i^m]).
\end{equation}
随后在前景区域内中心化并归一化：
\begin{align}
 \tilde z_i^{m\rightarrow n}&=z_i^{m\rightarrow n} -
 \frac{\sum_k w_k^m z_k^{m\rightarrow n}}{\sum_k w_k^m+\epsilon},\\
 g_i^{m\rightarrow n}&=\exp\{\beta\tanh(\tilde z_i^{m\rightarrow n})\},\\
 \hat g_i^{m\rightarrow n}&=\frac{g_i^{m\rightarrow n}}
 {\frac{\sum_k w_k^m g_k^{m\rightarrow n}}{\sum_k w_k^m}}.
\end{align}
因此前景区域内 $\hat g$ 的加权均值为 1，$\beta$ 控制 token 间重分配幅度。方向预算没有被放入这一步的均值归一化，避免全局方向项被代数抵消。

\subsection{跨模态 Value 更新与恒等初始化}
原生 MSHA 对拼接后的 RGB/TIR 模板 token 做联合注意力。将其输出按 query 模态和 key 来源模态分解，记 $A_{n\leftarrow m}$ 为接收模态 $n$ 的 query 到来源模态 $m$ 模板 token 的注意力子矩阵，$V_m$ 为对应 Value。DCVR-v2 只重加权跨模态项，同模态项保持原样：
\begin{equation}
 M_{n\leftarrow m}=A_{n\leftarrow m}
 \left[(b_{m\rightarrow n}\hat g_i^{m\rightarrow n})\odot V_{m,i}\right].
\end{equation}
所有同模态消息、Q/K 分数、原有残差和 CENTER 头保持不变。方向网络与 token 网络最后一层零初始化，使初始 $d=0$、$z=0$，从而 $b=1$、$\hat g=1$，候选输出与 parent 精确一致；零初始化投影仍可获得梯度。训练后由原始跟踪损失学习。零初始化、参数保存加载、梯度路径和数值稳定性是实现合同检查，不是性能证据。

\subsection{优化目标与限制}
初版不引入人工退化标签或额外质量监督，沿用基线的分类、框回归和质量损失：
\begin{equation}
 \mathcal{L}=\mathcal{L}_{\mathrm{base}}.
\end{equation}
因此路由是否与定位效用一致必须由独立诊断检验，不能从训练 loss 下降推断。该设计无法在两种模态都缺失目标信息时创造线索，也不处理重检测、长期记忆和目标出 crop 后恢复。实现还需记录路由极值、逐 token 权重、梯度、显存与延迟，防止均值统计掩盖少数 token 的异常放大。

\section{实验与分析}

\subsection{实验设置与评价指标}
受保护 parent 为 M1：\texttt{TBSI\_E\_MHSA\_FULL\_4090D\_seed42}，配置与 checkpoint 及其固定报告按实验资产数据库登记。其 LasHeR full-244 参考 AUC/PR/NPR/OP50/OP75 为 56.21/70.21/66.20/68.33/47.13，环境为 1$\times$ RTX 4090D、Python 3.8.20、PyTorch 1.9.0+cu111、CUDA 11.1、seed 42。该值是受保护记录，不用当前被改写的 M1 预测目录重新生成配对差值。

新候选必须有独立配置、输出目录、checkpoint 和资产 ID。训练/测试数据范围、manifest、预处理、更新策略、随机种子、评测脚本、权重和逐序列预测 hash 均须保存。主指标为 LasHeR AUC、PR、NPR、OP50、OP75；另报告逐序列配对差异、置信区间、胜负数、推理吞吐、参数量和延迟。full-244 单 seed 仅支持该次运行的描述性结论，不等于多 seed 稳健性。

表~\ref{tab:closure} 是拟议方法的验证清单，不是结果表。所有比较均须使用独立登记的资产；不满足数据、权重或评测合同的运行只能作为诊断或探索性结果。
\begin{table}[!t]
\centering
\caption{DCVR-v2 的最小证据闭环及当前状态}
\label{tab:closure}
\small
\begin{tabular}{@{}p{1.55cm}p{4.25cm}p{1.05cm}@{}}
\toprule
证据层 & 必需比较与读数 & 状态\\
\midrule
输入合同 & 正确 mask、训练/测试一致性、hash 与预处理 & 待完成\\
消息机制 & 方向预算、前景/背景分配、Value 改变量 & 待完成\\
局部定位 & CENTER rank、margin、offset/size、同 crop IoU & 待完成\\
受控退化 & 模态单独及联合退化、clean control & 待完成\\
训练贡献 & 同初值、同数据和优化预算的 on/off 配对训练 & 待完成\\
闭环结果 & LasHeR full-244、逐序列区间、失败与恢复统计 & 待完成\\
组件归因 & 前景、分布证据、方向预算、token 权重、等范数消融 & 待完成\\
成本复现 & FPS、延迟、显存、参数、环境及产物 hash & 待完成\\
\bottomrule
\end{tabular}
\end{table}

\subsection{基线参照与先导实验记录}
表~\ref{tab:existing} 汇总本文采用的性能参照、已完成的简化融合实验和先导模块结果。TBSI v1 一行由资产登记的 M1 相对 canonical raw 差值逐项反算；其余行是相应运行记录。各行用于交代研究起点和证据来源，不构成统一训练合同下的配对比较。
\begin{table*}[t]
\centering
\caption{性能参照与已有模块实验记录。TBSI v1 指标由资产记录差值算术倒推，并非独立复现；DCVR-v2 尚未评测。}
\label{tab:existing}
\small
\begin{tabular}{lccccccp{6.1cm}}
\toprule
运行/参照 & AUC & PR & NPR & OP50 & OP75 & FPS & 来源与说明\\
\midrule
TBSI v1 性能参照（算术倒推） & 54.82 & 68.53 & 64.59 & 66.52 & 45.80 & -- & 由 M1 相对 canonical raw 的登记差值逐项反算；不是独立重跑。\\
M1 简化融合（15 epoch） & 56.21 & 70.21 & 66.20 & 68.33 & 47.13 & NA & LasHeR full-244，单 seed；RTX 4090D，PyTorch 1.9.0+cu111；非 canonical 复现配置。\\
Early-conditioned MSHA gate & 56.56 & 70.58 & 66.76 & 68.69 & 47.64 & 53.14 & full-244 单 seed；联合训练且无 matched gate-off full15 control，旧配对分析来源链已失效。\\
DCVR-v1 full15 & 55.96 & 69.72 & 66.04 & 68.06 & 47.09 & 9.73 & full-244；训练 mask fallback 为 100\%，测试为 0\%；FPS 含诊断写盘，运行条件与 M1 不同。\\
DCVR-v2 & -- & -- & -- & -- & -- & -- & 仅为本文定义的后续候选；没有训练、推理或闭环结果。\\
\bottomrule
\end{tabular}
\end{table*}


\subsection{先导实验结果与证据边界}
表~\ref{tab:existing} 已在前文前置呈现。这里说明其解释范围：除受保护 M1 固定参考外，各候选均按自身登记运行报告，训练起点、比较控制和结果来源并不完全相同，不能据此计算统一的因果差值。

Early gate 的已登记 full-244 点估计均略高于 M1，但逐序列来源审计发现 parent 与 candidate 各有 7 条预测不匹配，当前 M1 目录重算 AUC 56.766，不符合受保护 M1 日志的 56.21。另一个 fixed-state A/B 中 Hann top-1 变化为 0/28，未观察到中心排名改善。故允许的结论是“曾记录方向性正点估计，但性能增益和机制未确立”。

DCVR-v1 的固定状态测试观察到路由非恒等，但 IoU 变化为 $-0.00062$，CENTER 排名仅 7 个事件中 1 个改变。受控退化 smoke 中，双模态 blur、contrast、lowlight 条件的 IoU 差分别为 $-0.00079$、$+0.00068$、$+0.00169$，clean 条件为 $+0.00554$；方向依赖条件，样本规模仅 8 个事件。quick3 候选相对控制的 AO 均值差为 $+0.04726$，主要由单一序列贡献，不能作为稳定增益证据。full15 DCVR-v1 汇总值可复算，但训练未获前景 mask、测试获得前景 mask；当前 M1 预测目录复算与 M1 保护日志不符，且 PyTorch/CUDA 版本不同。因此既不能把点估计下降单独归因于算子，也不能用现有配对 bootstrap 修复这些合同差异。

现有数据支持的判断限于：Early gate 与 DCVR-v1 能改变消息路径；DCVR-v1 full-244 已完成，汇总指标低于受保护 M1 的登记点估计。现有结果不支持两种候选已稳定改善 M1，也不支持路由信号已准确估计目标定位效用，或受控退化结果可推广至真实退化分布。DCVR-v1 结果受 mask 合同、信号与定位效用错位、跨软件结果来源及单 seed 影响的相对程度尚未区分；DCVR-v2 的各组成部分是否能依次改变消息分配、中心读出与闭环轨迹仍待检验。

\subsection{机制诊断与受控退化验证}
完整验证按阶段推进，前一阶段未通过时停止升级，不将局部现象包装成完整性能结论：
\begin{enumerate}
  \item \textbf{实现合同：} 检查 disabled parity、mask 训练/推理一致、张量形状、零初始化、保存加载、有限值、单 batch 梯度和冻结参数状态。
  \item \textbf{固定状态机制：} 在不与信号拟合重叠的 train-side 序列上，逐帧记录前景/背景 token 权重、双向消息预算、Value 改变量、CENTER 加窗前后排名、偏移/尺度解码和同 crop IoU。对照正确 mask、等面积移位 mask、均匀 token 权重及无方向路由。
  \item \textbf{受控退化：} 对 RGB 低照度/模糊、TIR 对比度/噪声及局部遮挡进行成对干预，检查退化来源方向是否受抑、另一方向是否获得相对预算，以及 clean 条件是否受损。合成干预只支持所测变换范围内的机制判断。
  \item \textbf{匹配代理训练：} 从相同 SOT 初始 state 克隆公共参数与新参数初值，使用相同 manifest、seed、优化器、步数和 worker 合同，只改变 DCVR-v2 开关，并与继续训练 control 成对比较。
  \item \textbf{闭环 full-244：} 仅在合同、固定状态和匹配代理满足预注册门槛后执行；固定权重、评测口径和 tracker 参数，保留逐序列结果、错误片段与恢复轨迹。
\end{enumerate}

\begin{figure*}[!t]
\centering
\fbox{\parbox[c][3.8cm][c]{0.94\textwidth}{\centering
\textbf{图件预留：固定状态机制诊断与受控退化结果}\par
正式图建议包含：(a) 两方向预算随 RGB/TIR 退化强度的变化；(b) 前景与背景 token 权重分布；(c) CENTER 排名或响应图变化；(d) clean 与退化条件下的同 crop 定位读数。\par
\textit{所有坐标、数值与可视化均待实验完成后填入；当前不构成结果。}}}
\caption{DCVR-v2 路由机制和受控退化诊断的预留版式。}
\label{fig:diagnostics}
\end{figure*}

\subsection{整体比较与消融实验版式}
若前述机制与代理训练达到预注册门槛，后续将在同一数据合同、评测器和软件环境下完成 full-244 比较。表~\ref{tab:main-result} 预留最终对比位置；当前候选结果尚不满足统一的匹配协议，故不填入先导点估计，也不据此作性能判断。
\begin{table*}[!t]
\centering
\caption{匹配协议下的 LasHeR full-244 整体比较（结果待实验完成后填写）}
\label{tab:main-result}
\small
\begin{tabular}{lcccccc}
\toprule
方法 & AUC$\uparrow$ & PR$\uparrow$ & NPR$\uparrow$ & OP50$\uparrow$ & OP75$\uparrow$ & FPS$\uparrow$\\
\midrule
受保护 M1（固定参考，非本轮配对重测） & -- & -- & -- & -- & -- & --\\
匹配 control（待运行） & -- & -- & -- & -- & -- & --\\
DCVR-v2（待运行） & -- & -- & -- & -- & -- & --\\
\bottomrule
\end{tabular}
\end{table*}

消融实验需要分别去除前景条件、匹配分布证据、跨模态一致性、方向预算和 token 级分配，并设置等范数缩放控制；同时报告新增参数、显存、延迟与纯推理 FPS。表~\ref{tab:ablation} 是待填版式，组件是否保留应由匹配实验决定。
\begin{table}[!t]
\centering
\caption{DCVR-v2 组件消融与计算成本（待实验）}
\label{tab:ablation}
\small
\begin{tabular}{@{}p{2.65cm}cccc@{}}
\toprule
设置 & AUC & CENTER & 参数量 & 延迟\\
\midrule
完整 DCVR-v2 & -- & -- & -- & --\\
去除前景条件 & -- & -- & -- & --\\
去除分布/一致性证据 & -- & -- & -- & --\\
去除方向预算 & -- & -- & -- & --\\
去除 token 分配 & -- & -- & -- & --\\
等范数缩放控制 & -- & -- & -- & --\\
\bottomrule
\end{tabular}
\end{table}

\begin{figure}[!t]
\centering
\fbox{\parbox[c][3.5cm][c]{0.92\columnwidth}{\centering
\textbf{图件预留：典型序列定性结果}\par
待补充相同帧、相同状态下的 RGB/TIR 输入、目标框、CENTER 响应与跟踪轨迹；需包含 clean、单模态退化及失败样例。\par
\textit{不使用示意图替代真实预测结果。}}}
\caption{代表性序列的定性比较（待补充真实输出）。}
\label{fig:qualitative}
\end{figure}

\subsection{讨论与局限}
DCVR-v2 的可检验价值在于把“消息来源的相对可靠性”拆成方向与 token 两个尺度，而不是给 gate 堆叠更多输入。熵低可能意味着目标匹配集中，也可能意味着背景纹理造成稳定错误峰；跨模态一致性高可能代表互补，也可能代表共同背景。因此路由可视化必须与 CENTER 响应和最终解码框一起检查。若消息权重变化没有传递到定位读出，就不能据机制图宣称方法有效。

当前 full-244 证据还有比较来源问题。M1 的受保护日志是唯一固定参照；若需要新候选的序列级因果比较，必须在独立输出目录中使用受保护 checkpoint 和候选 checkpoint、同一数据 manifest、同一评测器及同一软件环境重新生成结果，并验证每条序列预测 hash。不得覆盖 M1 记录或复用其输出目录。单 seed 结果只能报告为单次观察；资源预算不足以完成多 seed 时，应相应收窄结论。

此外，现有机制诊断指出的系统级风险链为“crop 内定位失配—反馈漂移—目标出界—恢复不足”。DCVR-v2 直接作用于模板 Value 分配，不直接修改选峰、框解码、状态提交或重检测。因此即使前两层证据成立，仍需展示其确实减少闭环失配或缩短失败持续时间；否则结论应限定为局部表示机制，不声称修复完整失败链。

\clearpage
\section{结论}
本文将 RGB-T 跟踪中的跨模态消息分配表述为前景条件下的方向预算与 token 级重分配问题，并给出 DCVR-v2 的可证伪设计。现有资产尚未验证该版本。已完成的 Early gate 与 DCVR-v1 实验提供了先导信号，也暴露出定位读出未改善、训练/测试 mask 不一致及结果来源不可严格配对等限制。当前最强结论是：跨模态 Value 可以被调制，但其目标效用和闭环收益尚未建立。只有通过合同一致的机制诊断、匹配训练对照和独立 full-244 闭环评测后，才能判断 DCVR-v2 是否构成有效跟踪改进。

\balance
\bibliographystyle{IEEEtran}
\bibliography{references}
\end{document}

